\section{Konsep Dasar Komunikasi}
\label{sec:konsep-komunikasi}

Komunikasi pada hakikatnya adalah pemindahan pesan dari satu pihak ke pihak lain
melalui suatu saluran. Dalam kajian kriptografi, komunikasi sengaja dimodelkan
secara sederhana: dua pihak yang berkepentingan, satu saluran, dan satu pihak lain
yang tidak diundang. Kesederhanaan model ini bukan kelemahan, melainkan justru yang
membuat persoalan keamanan informasi dapat dirumuskan secara tajam. Sebelum
membicarakan enkripsi, dekripsi, atau algoritma apa pun, kita perlu lebih dahulu
menyepakati siapa saja yang terlibat di dalamnya dan apa yang mungkin mereka
lakukan.

\subsection{Pihak yang Berkomunikasi dan Salurannya}

Dalam dunia nyata, komunikasi terjadi antara dua pihak, misalnya Alice (pengirim)
dan Bob (penerima). Mereka dapat berkomunikasi melalui berbagai saluran publik
seperti surat menyurat (pos), telepon, surel, serta pesan singkat dan aplikasi
perpesanan. Kata ``publik'' di sini penting dan perlu ditekankan sejak awal: saluran
tersebut tidak dirancang untuk menjaga kerahasiaan, sehingga apa pun yang melewatinya
harus dianggap dapat terlihat oleh pihak lain. Justru karena itulah kriptografi
diperlukan.

Pihak yang terlibat tidak selalu berupa manusia; bisa juga berupa interaksi antara
manusia dengan mesin (misalnya mesin penjawab atau layanan perbankan otomatis), atau
mesin dengan mesin (misalnya \textit{client} dan \textit{server} dalam jaringan
komputer). Pola terakhir ini kini mendominasi: sebagian besar komunikasi terenkripsi
di dunia modern terjadi tanpa manusia di salah satu ujungnya, misalnya ketika
peramban web mengambil halaman dari peladen, ketika aplikasi ponsel menyinkronkan
data, atau ketika dua perangkat Internet of Things bertukar pembacaan sensor.

Dalam pembahasan selanjutnya, demi kemudahan, pengirim pesan ditokohkan sebagai
Alice dan penerima pesan ditokohkan sebagai Bob. Konvensi penamaan tokoh ini lazim
dipakai dalam literatur kriptografi agar pembahasan terasa lebih manusiawi
ketimbang menggunakan simbol $A$, $B$, dan $C$.

\subsection{Penyusup dan Kemampuannya}

Tokoh ketiga dalam model ini adalah penyusup (\textit{intruder}), yaitu pihak yang
menyadap, mengintersepsi, menghapus, menambah, atau mengubah pesan. Sebutan lain
untuknya antara lain \textit{eavesdropper}, \textit{enemy}, \textit{adversary},
\textit{interceptor}, dan \textit{bad guy}; nama tokohnya bermacam-macam pula, seperti
Eve, Carol, Trudy, atau Mallory.

Salah satu perumusan paling ringkas mengenai mengapa keamanan komunikasi menjadi
persoalan sulit datang dari Ronald Rivest, salah satu pencipta algoritma RSA:
\textit{``cryptography is about communication in the presence of adversaries''}
--- kriptografi adalah tentang berkomunikasi ketika ada pihak yang berniat jahat.
Rumusan ini menegaskan bahwa asumsi adanya lawan bukan tambahan belaka, melainkan
bagian tak terpisahkan dari definisi masalahnya.

Pertanyaan yang wajar muncul kemudian: apa saja yang dapat dilakukan pihak ketiga?
Jawabannya: banyak. Menurut pembahasan pada \citet{stallings2017}, setidaknya ada
lima tindakan berikut.

\begin{enumerate}
    \item \textbf{Penyadapan (\textit{eavesdrop})}, yaitu mengintersepsi pesan yang
          melintas tanpa mengubahnya. Penyerang pasif seperti ini sulit dideteksi
          karena tidak meninggalkan jejak pada lalu lintas data.
    \item \textbf{Penyisipan pesan (\textit{insert})} ke dalam koneksi yang sedang
          berlangsung, sehingga pesan palsu tampak berasal dari pihak yang sah.
    \item \textbf{Penyamaran (\textit{impersonation})}, yaitu memalsukan alamat
          sumber pada paket data, atau bahkan sembarang ruas pada paket tersebut.
    \item \textbf{Pembajakan sesi (\textit{hijacking})}, yaitu mengambil alih koneksi
          yang sedang berjalan dengan menyingkirkan salah satu pihak lalu menempatkan
          diri di posisinya.
    \item \textbf{Serangan penolakan layanan (\textit{denial of service})}, yaitu
          mencegah layanan dipakai pihak lain, misalnya dengan membebani sumber daya
          sampai jenuh.
\end{enumerate}

Daftar ini memperlihatkan bahwa ancaman tidak berhenti pada pembacaan pesan.
Penyusup dapat pula mengubah, memalsukan, menambah, menghapus, bahkan melumpuhkan
komunikasi. Karena itu, satu mekanisme tunggal tidak akan pernah cukup; kriptografi
menyediakan sekumpulan layanan yang masing-masing menanggulangi sebagian ancaman.
Layanan-layanan itulah yang dibahas pada bagian berikutnya.
